\section{Reflection}

Below an inaccessible, we can ask for more than $\Sigma_1$-correctness: since $V_{\mu}$ is a set, full elementarity makes sense.

\subsection{Reflection Below an Inaccessible}

Here is the precise statement.

\begin{boxtheorem}
    Suppose $\mu$ is a strongly inaccessible cardinal. Let $C = \setst{\kappa < \mu}{V_{\kappa} \prec V_{\mu}}$. Then $C$ is a club subset of $\mu$.
\end{boxtheorem}

\begin{remark}
    $V_{\kappa} \prec V_{\mu}$ \textbf{is} a formula in the language of set theory because both $V_{\kappa}$ and $V_{\mu}$ are \underline{sets}. It would be a different story if they were classes.

    ``This is model theory. Model theory is math. And math is embedded in set theory. So this makes sense!''
\end{remark}

\begin{proof}
    Clearly $C$ is closed in $\mu$. If $C \cap \lambda \ne \emptyset$ and $\lambda = \sup(C \cap \lambda)$, then we need $\lambda \in C$ (meaning of closed).

    For such a $\lambda$, we have $V_{\lambda} = \bigcup_{\kappa < \lambda} V_{\kappa} = \bigcup_{\kappa \in C \cap \lambda} V_{\kappa}$.

    \begin{itemize}
        \item It follows, for $\kappa < \kappa' < \mu$ with $\kappa, \kappa' \in C$, that $V_{\kappa} \prec V_{\kappa'}$.

        \item Then, $\grpres{V_{\kappa}}{\kappa \in C \cap \lambda}$ is a chain of elementary substructures of $V_{\mu}$, so its union is also an elementary substructure of $V_{\mu}$.

        \item $V_{\lambda} \prec V_{\mu}$, so $\lambda \in C$.

        \item It also follows that $V_{\kappa} \prec V_{\lambda}$ for $\kappa \in C \cap \lambda$.
    \end{itemize}

    Now let $\kappa_0 < \mu$ be an arbitrary uncountable cardinal. By Downwards Löwenheim-Skolem, we have $X_0 \prec V_{\mu}$ with $V_{\kappa_0} \cup \set{\kappa_0} \ssq X_0$ and $\abs{X_0} = \abs{V_{\kappa_0} \cup \set{\kappa_0}} = \beth_{\kappa_0} < \beth_{\mu} = \mu$.

    % IDEA: Take X_0, fill in things of smaller rank, then take X_1, fill in evrything of smaller rank, and so on. Each X_i is obtained by applying DLS to $V_{\mu}$. We then get $V_{\kappa} \bigcup_{n < \omega} V_{\kappa_n} = \bigcup_{n < \omega} X_n \prec V_{\kappa}$ where we konw the chain is elementary by the elementary chain theorem.
    \sorry
\end{proof}

\subsection{Hierarchies}

To state reflection in general, we give a name to the shape that $\grpres{V_{\alpha}}{\alpha \in \OR}$ and the $H_{\kappa}$ share.

\begin{boxdefinition}[Hierarchy]
    We say that a class function $\grpres{Z_{\alpha}}{\alpha \in C}$ is a \textbf{hierarchy} of sets iff
    \begin{itemize}
        \item each $Z_{\alpha}$ is a transitive set
        \item $C$ is a club class of ordinals
        \item for every $\alpha < \beta$ from $C$, $Z_{\alpha} \subset Z_{\beta}$
        \item for every $\beta$ a limit point\footnote{as $C$ is club, such $\beta$s do also lie in $C$} of $C$,
        \begin{align*}
            Z_{\beta} = \bigcup_{\alpha \in C \cap \beta} Z_{\alpha}
        \end{align*}
    \end{itemize}
    We usually write
    \begin{align*}
        Z = \bigcup_{\alpha \in C} Z_{\alpha}
        = \setst{x}{\exists \alpha \in C \st x \in Z_{\alpha}}
    \end{align*}
    (unless, of course, we already have a name for our hierarchy).
\end{boxdefinition}

We have already seen two of these.

\begin{boxexample}
    The following are all examples of hierarchies we've seen:
    \begin{align*}
        \grpres{V_{\alpha}}{\alpha \in \OR}
        \qquad , \qquad
        \grpres{H_{\kappa}}{\kappa \text{ an infinite cardinal}}
    \end{align*}
\end{boxexample}

Indeed, there will be examples where $Z \ne V$.

\subsection{The Reflection Theorem}

Every hierarchy reflects any finitely many formulae down to club many of its levels.

\begin{boxtheorem}[Reflection theorem scheme, proved in $\ZF$]
    Suppose $\varphi_i$ for $i < \ell < \omega$ are formulae in the language of set theory. Suppose that $\grpres{Z_{\alpha}}{\alpha \in C}$ is a hierarchy. Then there's a club class $D \ssq C$ such that for every $\alpha \in D$ and $i < \ell$, $\vp_i$ is absolute between $Z_{\alpha}$ and $Z$.
\end{boxtheorem}
\begin{proof}
    WLOG, assume that for all $j < \ell$ and proper subformulae $\psi$ of $\varphi_j$, there is some $i < j$ such that $\psi$ is $\vp_{i}$. Ie, we are setting up an induction on the complexity of $\vp$ by setting up an induction on ordinals.

    Let's define a class sequence of ordinals $\grpres{\alpha_{\eta}}{\eta \in \OR}$ as follows.
    \begin{itemize}
        % [CORRECTED] was: $\alpha_0 = 0$ --- $Z_{\alpha}$ is only defined for $\alpha \in C$, and $0$ need not be in $C$ (it is not for $\grpres{H_{\kappa}}{\kappa \text{ an infinite cardinal}}$), while the successor step uses $Z_{\alpha_0}$; the board also says to maintain $\alpha_{\eta} \in C$ for all $\eta$
        \item \underline{Base:} Define $\alpha_0 = \min(C)$.
        \item \underline{Successor:} Given $\alpha_{\eta}$, define $\alpha_{\eta + 1}$ to be the least $\beta \in C \setminus \parenth{\alpha_{\eta} + 1}$ so that for all $j < \ell$, if $\vp_{j}\of{\bar{u}}$ is of the form $\exists v, \vp_i\of{\bar{u}, v}$, then for all $\bar{x}$ from $Z_{\alpha_{\eta}}$ that match $\bar{u}$, if there is $y \in Z$ such that $\vp_i\of{\bar{x}, y}^Z$, then there is $y \in Z_{\beta}$ such that $\vp_i\of{\bar{x}, y}^Z$.
        \item \underline{Limit case:} If $\theta$ is a limit ordinal, define
        \begin{align*}
            \alpha_{\theta} = \sup_{\eta < \theta} \of{\alpha_{\eta}} \in \lim(C) \subset C
        \end{align*}
        (because $C$ is club).
    \end{itemize}

    The successor step is where the work happens. A tuple $\bar{x}$ from $Z_{\alpha_{\eta}}$ may have a witness $y'$ only higher up, in some $Z_{\beta'}$, and $\alpha_{\eta + 1}$ is chosen high enough to catch a witness for every such $\bar{x}$ at once:
    \begin{figure}[H]
        \centering
        \begin{tikzpicture}[dot/.style = {circle, fill, inner sep = 1.5pt}, every node/.style = {font = \small}]
            \drawcone{2.2}{5}{$Z$}
            \drawconelevel{2.2}{5}{1.6}{$Z_{\alpha_{\eta}}$}
            \drawconelevel{2.2}{5}{4.2}{$Z_{\beta'}$}
            \node[dot, label = {right:$\bar{x}$}] at (-0.15, 1.0) {};
            \node[dot, label = {right:$y'$}] at (0.3, 3.6) {};
        \end{tikzpicture}
    \end{figure}
    Only existential formulae need this treatment. We may take every $\vp_j$ to be built from atomic formulae using $\neg$, $\land$ and $\exists$ alone, writing $\forall$ as $\neg \exists \neg$. Relativization leaves atomic formulae unchanged, so they mean the same in $Z_{\alpha}$ as in $Z$, and it commutes with $\neg$ and $\land$, so absoluteness passes through them. The one place it can fail is $\exists v$, where $Z$ may have a witness that $Z_{\alpha}$ lacks.

    The successor step makes sense. For each $j < \ell$ of the form above and each $\bar{x}$ from $Z_{\alpha_{\eta}}$ that has a witness $y \in Z$, there is a least $\gamma \in C$ with a witness in $Z_{\gamma}$, as $Z = \bigcup_{\gamma \in C} Z_{\gamma}$. There are only set many such pairs $\parenth{j, \bar{x}}$, as $Z_{\alpha_{\eta}}$ is a set, so these $\gamma$ have a supremum, and as $C$ is unbounded there is some $\beta \in C$ above it and above $\alpha_{\eta}$. Any such $\beta$ works, since $Z_{\gamma} \ssq Z_{\beta}$ for $\gamma \leq \beta$ in $C$.

    Then take $D$ to be $\setst{\alpha_{\theta}}{\theta \text{ is a limit ordinal}}$. The sequence $\grpres{\alpha_{\eta}}{\eta \in \OR}$ is strictly increasing and continuous at limits, so $D$ is a club class: it is unbounded since $\alpha_{\theta} \geq \theta$, and closed since a supremum of $\alpha_{\theta}$s over limit ordinals $\theta$ in some set $\Theta$ with no largest element is $\alpha_{\sup \Theta}$, and $\sup \Theta$ is a limit ordinal. And $D \ssq C$ by the limit case.

    Finally, fix $\alpha = \alpha_{\theta} \in D$; we show that each $\vp_i$ is absolute between $Z_{\alpha}$ and $Z$, by induction on $i$. The atomic and propositional cases are as above, so suppose $\vp_j\of{\bar{u}}$ is $\exists v \, \vp_i\of{\bar{u}, v}$ with $i < j$, and let $\bar{x}$ from $Z_{\alpha}$ match $\bar{u}$.
    \begin{description}
        \item[\underline{From $Z_{\alpha}$ to $Z$.}] If $y \in Z_{\alpha}$ and $\vp_i\of{\bar{x}, y}^{Z_{\alpha}}$, then $y \in Z$ and $\vp_i\of{\bar{x}, y}^Z$ by the induction hypothesis, so $\vp_j\of{\bar{x}}^Z$.

        \item[\underline{From $Z$ to $Z_{\alpha}$.}] Suppose $\vp_j\of{\bar{x}}^Z$. As $\alpha_{\theta}$ is a limit point of $C$, $Z_{\alpha_{\theta}} = \bigcup_{\gamma \in C \cap \alpha_{\theta}} Z_{\gamma} = \bigcup_{\eta < \theta} Z_{\alpha_{\eta}}$, and $\bar{x}$ is a finite tuple, so $\bar{x}$ is from $Z_{\alpha_{\eta}}$ for some $\eta < \theta$. By the choice of $\alpha_{\eta + 1}$, there is $y \in Z_{\alpha_{\eta + 1}} \ssq Z_{\alpha}$ with $\vp_i\of{\bar{x}, y}^Z$, and by the induction hypothesis $\vp_i\of{\bar{x}, y}^{Z_{\alpha}}$. So $\vp_j\of{\bar{x}}^{Z_{\alpha}}$.
    \end{description}
\end{proof}

Here's an interesting corollary.

Let $\sigma_0, \ldots, \sigma_{\ell}$ be finitely many axioms of $\ZFC$. By the reflection theorem scheme, there are club many $\alpha \in \OR$ so that $\forall  i \leq \ell$, $\sigma_i^{V_{\alpha}}$. Equivalently, $\forall i \leq \ell$, $V_{\alpha} \models \sigma_i$. Rephrasing:

\begin{boxcorollary}
    For every finite sub-theory $T$ of $\ZFC$,
    \begin{align*}
        \ZFC \vdash \text{``$T$ has a transitive model''}
    \end{align*}
    where the quotes (``read with a different tone of voice'')
\end{boxcorollary}

With Choice, we can make the model countable as well. Given $\alpha$ with
\begin{align*}
    V_{\alpha} \models \bigwedge_{i \leq \ell} \sigma_i
\end{align*}
we can apply Downwards Löwenheim-Skolem to get some countable $X \prec V_{\alpha}$, so that $X \models \bigwedge_{i \leq \ell} \sigma_i$ too. As $V_{\alpha}$ is transitive, it satisfies Extensionality, and hence so does $X$; so $\in$ is extensional on $X$, as well as well-founded and set-like. Then apply Mostowski (\Cref{Ch2:Thm:Mostowski}) to get $X \simeq M$ with $M$ transitive and countable, and $M \models \bigwedge_{i \leq \ell} \sigma_i$, because satisfaction is preserved by isomorphisms. So in fact $\ZFC$ proves that every finite sub-theory of $\ZFC$ has a countable transitive model.
